\subsection{Formal effectivity of the quadratic modulus}
\label{subsec:formal-effectivity-q-modulus}

The preceding formal source-triviality result can be sharpened in two
independent directions.  First, the framed root-translation equation has an
exact annihilator and degree law over every coefficient ring.  Second, the
resulting compatible Artin isomorphisms do not algebraize to a stable
polynomial left--right equivalence over \(\C[[s]]\).  Thus the stable
left--right groupoid itself fails formal effectivity at the quadratic
modulus, even though the quantitative lower bound below is proved only in the
framed translation groupoid.  The complete stable \(q\)-classification is
the load-bearing input for generic-fiber nonexistence; the existing all-order
coefficientwise formal source triviality is overlapping background sharpened
by the explicit calculation below.  The exact annihilator law,
nonalgebraizability, unrestricted complexity escape, and diagonal obstruction
are the new deductions in this subsection.

For a commutative \(\Q\)-algebra \(R\), let
\[
 A_\alpha(c)=c(1+\alpha c),\qquad
 B_{\alpha,q}(c)=-2-4\alpha c+q\alpha^2c^2,
\]
and write \(F_{\alpha,q}=G_{A_\alpha,B_{\alpha,q}}\).  Put
\(\delta=q'-q\).

For \(\phi(c)\in cR[c]\), define
\[
 \Theta_\phi(x,y,z)=
 \left(x,y+\phi(c),z-3\frac{\phi(c)}x\right),
\]
\[
 \ell_\phi=3A\phi^2+2B\phi,
 \qquad
 \eta_\phi=A\phi^3+B\phi^2,
\]
and
\[
 \Xi_\phi(a,b,c)=
 \left(a-\frac12\phi(c)b-\frac12\eta_\phi(c),
       b+\ell_\phi(c),c\right).
\]
The source map is polynomial because \(c/x=2-3xy-x^2z\), fixes \(c\),
and shifts \(t\) by \(\phi(c)\).  Direct expansion gives
\begin{equation}
\label{eq:q-pairwise-root-translation}
 G_{A,B+3A\phi}=\Xi_\phi\circ G_{A,B}\circ\Theta_\phi.
\end{equation}

\begin{theorem}[Exact framed effectivity law]
\label{thm:q-framed-effectivity-law}
Let \(D\ge0\).  A \(c\)-fixed framed root translation of
\(c\)-degree at most \(D\) carries \(F_{\alpha,q}\) to
\(F_{\alpha,q'}\) if and only if
\[
 \delta\alpha^{D+2}=0.
\]
When it exists, it is unique and is given by
\[
 \phi_D(c)=\frac\delta3\alpha^2c
 \sum_{j=0}^{D-1}(-\alpha c)^j,
\]
where the sum is empty for \(D=0\).  Before imposing the annihilator
condition, the residual is exactly
\[
 B_{\alpha,q'}-B_{\alpha,q}-3A_\alpha\phi_D
 =(-1)^D\delta\alpha^{D+2}c^{D+2}.
\]
Consequently, if
\[
 N=\min\{n\ge2:\delta\alpha^n=0\}
\]
exists, then \(N=2\) means that the frames already agree and the
translation is zero; for \(N\ge3\), its exact degree is \(N-2\).  If no
such \(N\) exists, there is no polynomial framed translation.
\end{theorem}

\begin{proof}
The coefficient equation is
\[
 3c(1+\alpha c)\phi(c)=\delta\alpha^2c^2.
\]
Canceling \(c\) and writing
\(\phi=p_1c+\cdots+p_Dc^D\) gives
\[
 3p_1=\delta\alpha^2,\qquad
 p_i=-\alpha p_{i-1}\ (2\le i\le D),\qquad
 \alpha p_D=0.
\]
Thus
\[
 p_i=\frac\delta3(-1)^{i-1}\alpha^{i+1},
\]
and the terminal equation is precisely
\(\delta\alpha^{D+2}=0\).  This proves existence and uniqueness.  The
finite geometric-series identity gives the displayed residual.  For \(N\ge3\), minimality makes the coefficient of \(c^{N-2}\)
nonzero, proving the exact degree.  When \(N=2\), the frames already
coincide and \(\phi=0\).
\end{proof}

\begin{proposition}[Exact framed orbit cokernel]
\label{prop:q-framed-orbit-cokernel}
Let
\[
 \mu_\alpha:cR[c]\longrightarrow c^2R[c],
 \qquad \phi\longmapsto3A_\alpha\phi.
\]
Then
\[
 \operatorname{coker}(\mu_\alpha)
 \simeq R[c]/(1+\alpha c),
\]
and the difference between the \(q'\)- and \(q\)-frames is the class
\[
 \frac{q'-q}{3}\alpha^2\bmod(1+\alpha c).
\]
For \(R=\C[[s]]\) and \(\alpha=s\), this cokernel is
\[
 \C[[s]][c]/(1+sc)\simeq\C((s)).
\]
It is nonzero, but its reduction modulo every \(s^M\), and hence its
\(s\)-adic completion, is zero.
\end{proposition}

\begin{proof}
Divide the source by \(c\), the target by \(c^2\), and the map by the unit
three.  The resulting map is multiplication by \(1+\alpha c\).  In the
special case, the quotient relation is \(c=-s^{-1}\), giving
\(\C[[s]][1/s]=\C((s))\).  Since \(s\) is invertible there, all
\(s\)-power quotients vanish.
\end{proof}

\begin{corollary}[Ramification and obstruction staircase]
\label{cor:q-ramification-degree-law}
Let \(R_M=\C[s]/(s^M)\), let
\(\alpha=s^eu(s)\) with \(u(0)\ne0\), and let \(q\ne q'\).  The exact
framed translation degree is
\[
 D_M=\max\left(0,\left\lceil\frac Me\right\rceil-2\right).
\]
For the unramified arc \(\alpha=s\), this is \(D_M=M-2\) for \(M\ge3\),
and the optimal degree-\(D\) residual is
\[
 (-1)^D(q'-q)s^{D+2}c^{D+2}.
\]
For \(D\ge1\), the canonical source and target automorphisms in
\eqref{eq:q-pairwise-root-translation} have exact ordinary degrees
\[
 \deg\Theta_{\phi_D}=4D,
 \qquad
 \deg\Xi_{\phi_D}=D+1.
\]
\end{corollary}

\begin{proof}
The nilpotence index of \(s^eu(s)\) in \(R_M\) is
\(\lceil M/e\rceil\), so the first assertion follows from
\cref{thm:q-framed-effectivity-law}.  The residual formula is its displayed
identity with \(\alpha=s\).  Since \(c(x,y,z)\) has degree four and the top
coefficient of \(\phi_D\) is nonzero, the source degree is \(4D\).  On the
target the term \(\phi_D(c)b\) has degree \(D+1\), while nilpotence and the
explicit formulas for \(\ell_\phi,\eta_\phi\) bound all remaining
corrections by \(D\).
\end{proof}

The same degree survives the residual affine transformations of the
normalized conductor chart.  Indeed, over a local ring with \(\alpha\) in
the maximal ideal, suppose
\[
 C=uc+v,\qquad T=\nu t+h(c)
\]
satisfies
\[
 A_\alpha(C)\nu=\kappa A_\alpha(c),
\qquad
 B_{\alpha,q'}(C)+3A_\alpha(C)h(c)=
 \kappa B_{\alpha,q}(c).
\]
Coefficient comparison gives
\[
 v=0,\quad \kappa=1,\quad \nu u=1,\quad
 \alpha(u-1)=0,
\]
and then
\[
 h(c)=\frac{q-q'}{3u}\frac{\alpha^2c}{1+\alpha c}.
\]
Multiplication by the unit \(u^{-1}\) does not alter the exact degree in
\cref{thm:q-framed-effectivity-law}.

\begin{theorem}[Failure of formal effectivity in the stable groupoid]
\label{thm:q-stable-formal-noneffectivity}
Let
\[
 \mathcal F_q=F_{s,q}
\]
be regarded as a polynomial Keller map over \(\C[[s]]\).  If \(q\ne q'\),
then:
\begin{enumerate}[label=(\roman*)]
\item for every \(M\ge1\), the reductions of
\(\mathcal F_q\) and \(\mathcal F_{q'}\) modulo \(s^M\) are ordinarily
polynomially left--right equivalent;
\item these equivalences may be chosen compatibly in \(M\); but
\item the two maps are not stably polynomially left--right equivalent over
\(\C[[s]]\).
\end{enumerate}
Consequently the stable isomorphism functor is not formally effective at
this pair of arcs.
\end{theorem}

\begin{proof}
For \(M\le2\) the frames agree.  For \(M\ge3\), use
\[
 \phi_M(c)=\frac{q'-q}{3}s^2c
 \sum_{j=0}^{M-3}(-sc)^j
\]
in \eqref{eq:q-pairwise-root-translation}.  The extra term in
\(\phi_{M+1}\) is divisible by \(s^M\), so the equivalences are compatible.

Suppose a stable polynomial equivalence existed over \(\C[[s]]\).  Base
change to an algebraic closure \(L\) of \(\C((s))\).  The diagonal
cubic-frame scaling normalizes the nonzero coefficient \(\alpha=s\) to
\(\alpha=1\), so the two generic fibers become \(G_q\) and \(G_{q'}\)
over \(L\).

Fix the stabilization dimension and the finite degrees of the polynomial
automorphisms and their inverses in this alleged equivalence.  Their
coefficients form an \(L\)-point of a finite-type affine \(\C\)-scheme cut
out by the inverse-composition equations and the left--right identity.  Since
\(\C\) is algebraically closed, nonemptiness gives a \(\C\)-point.  That
would be a stable polynomial equivalence between \(G_q\) and \(G_{q'}\)
over \(\C\), contradicting the complete \(q\)-classification.
\end{proof}

\subsection{Effective unframed complexity lower bounds}
\label{subsec:q-effective-unframed-bound}

The preceding framed calculation has a linear degree law.  We now give a
weaker but completely unframed lower bound.  No recovery of the cubic frame,
conductor chart, or escaping section is assumed.

Fix distinct \(q,q'\in\C\), put
\[
 F_q=G_{c(1+sc),\,-2-4sc+qs^2c^2},
 \qquad R_M=\C[s]/(s^M),
\]
and fix a stabilization dimension \(m\ge0\) and a degree bound \(b\ge1\).
Set
\[
 n=3+m,\qquad
 T(n,b)=\binom{n+b}{n},\qquad
 N(n,b)=4nT(n,b),
\]
\[
 d_b=\max\{b+1,11\},\qquad h_b=2b.
\]

Introduce coefficient variables for four polynomial maps
\[
 \Phi,\Phi^{-1},\Psi,\Psi^{-1}:\A^n\longrightarrow\A^n
\]
of degree at most \(b\).  Their two-sided inverse equations and the stable
left--right identity
\[
 (F_{q'}\times\id_{\A^m})\circ\Phi
 =\Psi\circ(F_q\times\id_{\A^m})
\]
define an affine scheme \(E_{m,b}\) over \(\C[s]\).  It has
\(N=N(n,b)\) coefficient variables.  Every defining equation \(f_i\) obeys
\[
 \deg_X(f_i)\le d_b,
 \qquad
 \deg_s(f_i)\le h_b.
\]
Indeed, a composition-inverse coefficient has degree at most \(b+1\) in the
unknown coefficients; substitution into the degree-eleven map has coefficient
degree at most eleven; and a degree-\(b\) monomial in coordinates whose
coefficients have \(s\)-degree at most two has \(s\)-degree at most \(2b\).

\begin{lemma}[Reduction to \(N+1\) constant combinations]
\label{lem:q-N-plus-one-combinations}
Let \(k\) be an infinite field and
\(f_1,\ldots,f_r\in k[s,X_1,\ldots,X_N]\) have no common zero over
\(\overline{k(s)}\).  There are constants \(\lambda_{ji}\in k\),
\(0\le j\le N\), such that
\[
 g_j=\sum_i\lambda_{ji}f_i
\]
are nonconstant in the \(X\)-variables and have no common zero over
\(\overline{k(s)}\).
\end{lemma}

\begin{proof}
Over \(K=\overline{k(s)}\), let \(\Lambda=(\lambda_{ji})\) range over
\(\A_K^{(N+1)r}\), and consider the incidence scheme
\[
 I=\left\{(x,\Lambda):
 \sum_i\lambda_{ji}f_i(x)=0\text{ for }0\le j\le N\right\}.
\]
For each \(x\), the vector \((f_1(x),\ldots,f_r(x))\) is nonzero.  Each row
of \(\Lambda\) therefore satisfies one nontrivial linear equation, and
\[
 \dim I\le N+(N+1)(r-1)=(N+1)r-1.
\]
The closure of its projection to the coefficient space is proper.  Requiring
a row not to produce an element of \(k[s]\) removes only a proper linear
subspace.  Hence a nonempty open set of tuples works over \(k(s)\).

This open contains a tuple with entries in \(k\).  Choose a nonzero polynomial
over \(k(s)\) vanishing on the bad closed set and clear denominators.  A
nonzero polynomial in \(k[s,\Lambda]\) cannot vanish at every constant point
of \(k^{(N+1)r}\), since \(k\) is infinite.
\end{proof}

The generic fiber of \(E_{m,b}\) is empty.  Otherwise, after extension to an
algebraic closure of \(\C(s)\), diagonal scaling would normalize \(s\ne0\)
and give a stable equivalence between \(G_q\) and \(G_{q'}\).  With \(m\)
and \(b\) fixed, such an equivalence is a point of a finite-type scheme over
\(\C\); a point over an extension field would yield a complex point,
contradicting the complete stable \(q\)-classification.

Apply \cref{lem:q-N-plus-one-combinations} to obtain \(N+1\) polynomials
\(g_0,\ldots,g_N\) with no generic common zero and with the same degree
bounds.  The parametric effective Nullstellensatz
\cite[Theorem~0.5]{dAndreaKrickSombra2013} supplies
\(0\ne\alpha_{m,b}(s)\in\C[s]\) and \(a_j\in\C[s,X]\) such that
\[
 \alpha_{m,b}(s)=\sum_{j=0}^{N}a_jg_j
\]
and
\[
 \deg_s\alpha_{m,b}
 \le\sum_{\ell=0}^{N}
 \left(\prod_{j\ne\ell}\deg_Xg_j\right)\deg_sg_\ell
 \le (N+1)d_b^Nh_b.
\]

\begin{theorem}[Effective unframed truncation bound]
\label{thm:q-effective-unframed-truncation-bound}
If the two reductions modulo \(s^M\) admit a stable polynomial left--right
equivalence with exactly \(m\) stabilization variables and with
\[
 \deg\Phi,\deg\Phi^{-1},\deg\Psi,\deg\Psi^{-1}\le b,
\]
then
\[
 \boxed{
 M\le H(m,b):=
 2b\bigl(N(n,b)+1\bigr)d_b^{N(n,b)}.}
\]
\end{theorem}

\begin{proof}
An \(R_M\)-point of \(E_{m,b}\) annihilates every \(g_j\).  Evaluation of
the Bezout identity gives \(\alpha_{m,b}(s)=0\) in \(R_M\).  Since
\(\alpha_{m,b}\ne0\),
\[
 M\le\ord_s\alpha_{m,b}
 \le\deg_s\alpha_{m,b}
 \le2b(N+1)d_b^N.
\]
\end{proof}

\begin{corollary}[Fixed-stabilization degree rate]
\label{cor:q-fixed-stabilization-degree-rate}
Let \(b_{M,m}\) be the least common degree bound for an equivalence and its
four automorphism maps using exactly \(m\) added variables, and put
\(n=3+m\).  Then
\[
 \liminf_{M\to\infty}
 \frac{b_{M,m}}{(\log M/\log\log M)^{1/n}}
 \ge\left(\frac{n!}{4}\right)^{1/n}.
\]
In particular,
\[
 b_{M,0}\ge
 \left(\frac32\frac{\log M}{\log\log M}\right)^{1/3}(1-o(1)).
\]
\end{corollary}

\begin{proof}
For fixed \(n\),
\[
 N(n,b)=4n\binom{n+b}{n}
 =\frac4{(n-1)!}b^n+O_n(b^{n-1}),
\]
so
\[
 \log H(m,b)
 =\frac4{(n-1)!}b^n\log b+O_n(b^n).
\]
Asymptotic inversion gives the result.
\end{proof}

\begin{corollary}[Stabilization--degree tradeoff]
\label{cor:q-stabilization-degree-tradeoff}
Every such equivalence satisfies
\[
 M\le
 2b\bigl(32(m+3)2^{m+b}+1\bigr)
 (b+11)^{32(m+3)2^{m+b}},
\]
and consequently
\[
 m+b\ge\frac{\log\log M}{\log2}
       -O(\log\log\log M).
\]
\end{corollary}

\begin{proof}
Use
\(\binom{n+b}{n}\le2^{n+b}=2^{m+b+3}\), hence
\(N(n,b)\le32(m+3)2^{m+b}\), and substitute in
\cref{thm:q-effective-unframed-truncation-bound}.
\end{proof}

\begin{corollary}[Explicit unrestricted complexity rate]
\label{cor:q-unrestricted-complexity-rate}
Let
\[
 \kappa_M(q,q')=
 \min\max\{m,\deg\Phi,\deg\Phi^{-1},
              \deg\Psi,\deg\Psi^{-1}\}
\]
over all stable polynomial equivalences modulo \(s^M\).  Then
\[
 \boxed{
 \kappa_M(q,q')
 \ge\frac{\log\log M}{\log4}
      -O(\log\log\log M),}
\]
or equivalently
\[
 \boxed{
 \liminf_{M\to\infty}
 \frac{\kappa_M(q,q')}{\log\log M}\ge\frac1{\log4}.}
\]
More explicitly, \(\kappa_M(q,q')\le B\) implies
\[
 M\le
 2B\bigl(32(B+3)4^B+1\bigr)
 (B+11)^{32(B+3)4^B}.
\]
\end{corollary}

The result is deliberately worst-case: it treats the coefficients of four
bounded polynomial maps as independent variables and applies general
elimination.  It nevertheless proves an explicit rate for arbitrary stable
equivalences, with no frame-recovery hypothesis.  The sharp linear lower
bound remains a geometric problem.  The explicit framed construction gives
\(\kappa_M(q,q')\le4M-8\); matching this requires an intrinsic Artin-base
recovery theorem for the escaping-boundary chart.

The same non-effectivity argument works over an integral
\(\C\)-algebra \(R\) that is complete and separated for the
\(\alpha\)-adic topology, with \(\alpha\) a nonzero nonunit and
\(q,q'\in\C\) distinct.  Strictness of the powers of \(\alpha\) gives exact
framed degree \(M-2\) modulo \(\alpha^M\), and a complete-base stable
equivalence would contradict the complex generic-fiber classification after
finite-type descent.

The compatible framed translations have the unique limit
\[
 \widehat\phi(c)=\frac{q'-q}{3}\frac{s^2c}{1+sc}
 =\frac{q'-q}{3}\sum_{j\ge0}(-1)^js^{j+2}c^{j+1}.
\]
It lies in \(c\C[c][[s]]\), the \(s\)-adic completion of
\(c\C[s,c]\), but not in \(c\C[[s]][c]\).  Equivalently, if
\(\mathcal I_D(M)\) denotes framed isomorphisms of translation degree at
most \(D\), then
\[
 \varprojlim_M\varinjlim_D\mathcal I_D(M)\ne\varnothing,
 \qquad
 \varinjlim_D\varprojlim_M\mathcal I_D(M)=\varnothing.
\]
Thus completion and the bounded-degree filtration do not commute.

\begin{corollary}[Obstruction to an affine finite-presentation diagonal]
\label{cor:q-no-affine-fp-moduli-diagonal}
No algebraic stack can represent this stable polynomial left--right
groupoid near the two arcs, with its exact isomorphism notion, and have an
affine diagonal locally of finite presentation.
\end{corollary}

\begin{proof}
If such a stack existed, the isomorphism space between the two
\(\C[[s]]\)-objects would be affine of finite presentation, say
\(\Spec A\).  Since \(A\) is finitely presented and
\(\C[[s]]=\varprojlim_M\C[s]/(s^M)\), one has
\[
 \Hom(A,\C[[s]])
 \simeq
 \varprojlim_M\Hom(A,\C[s]/(s^M)).
\]
The compatible Artin isomorphisms would therefore algebraize, contradicting
\cref{thm:q-stable-formal-noneffectivity}.
\end{proof}
